\chapter[Fondements et revue de littérature]{Fondements théoriques et revue de littérature}
\label{ch:etat}

Ce chapitre donne au lecteur ce qui est nécessaire pour suivre la suite sans autre source : le schéma de signature MAYO et ses algorithmes (section~\ref{sec:schema}), leur traduction dans le code de référence et la place des noyaux étudiés (section~\ref{sec:code}), la représentation des données et l'arithmétique par \vpshufb{} sur $\mathbb{F}_{16}$ (section~\ref{sec:arith}), un panorama des analyses menées (section~\ref{sec:panorama}), puis les implémentations existantes et le positionnement de l'optimisation étudiée (sections~\ref{sec:revue} et suivantes).

\section{Fondements théoriques}
\subsection{Le schéma MAYO}\label{sec:schema}
\subsubsection{Des signatures multivariées à Oil-and-Vinegar, puis à UOV}
Une signature « multivariée » repose sur une fonction publique $\mathcal{P}:\mathbb{F}_q^n\to\mathbb{F}_q^m$ formée de $m$ polynômes de degré~2 en $n$ variables. Signer un message $M$, c'est trouver un point $s$ tel que $\mathcal{P}(s)=H(M)$, où $H$ est une fonction de hachage (avec un « sel » aléatoire) qui donne $m$ valeurs de $\mathbb{F}_q$. Vérifier une signature, c'est simplement calculer $\mathcal{P}(s)$. Le schéma Oil-and-Vinegar (OV), proposé par Patarin en 1997~\cite{patarin}, cache un secret : un sous-espace $\mathcal{O}\subset\mathbb{F}_q^n$ (« l'huile »), de dimension $o$, sur lequel $\mathcal{P}$ vaut zéro. Les $n$ variables se séparent en $v=n-o$ variables « vinaigre » et $o$ variables « huile ». Dans OV, on prend $o=m$. Pour signer, on choisit les variables vinaigre au hasard. Il reste alors un \emph{système linéaire} de $m$ équations avec $o$ inconnues (les variables huile), parce que, dans la base secrète, aucun terme ne multiplie deux variables huile entre elles. On le résout directement.

Dans la version d'origine, dite \emph{équilibrée}, on avait $v=o$. Cette version a été cassée en 1998 par Kipnis et Shamir~\cite{kipnisshamir}. La correction, proposée en 1999 par Kipnis, Patarin et Goubin, s'appelle \emph{Unbalanced Oil and Vinegar} (UOV)~\cite{uov} : on prend $v>o$ (beaucoup plus de vinaigre que d'huile) et on garde $o=m$. La signature se déroule comme dans OV (vinaigre au hasard, puis système linéaire de $m$ équations). Mais l'attaque contre la version équilibrée ne marche plus, et UOV n'a pas été cassé depuis. Son défaut est la taille de la clé publique : $n=v+o$ est grand, et la clé contient $m$ matrices triangulaires de taille $n\times n$.

\subsubsection{Le schéma MAYO}
MAYO~\cite{mayo,nibbling} part d'UOV (où $o=m$) et choisit un espace d'huile \emph{trop petit} ($o<m$). Cela réduit beaucoup la taille des clés, mais on ne peut plus résoudre le système directement : il y a $m$ équations pour seulement $o$ inconnues. La solution est le « fouettage » (\emph{whipping}). On étend la fonction publique $\mathcal{P}$ en une fonction $\mathcal{P}^\star:\mathbb{F}_q^{kn}\to\mathbb{F}_q^m$ qui prend $k$ blocs de $n$ variables :
\[\mathcal{P}^\star(x_1,\dots,x_k)=\sum_{i=1}^{k}E_{ii}\,\mathcal{P}(x_i)+\sum_{i=1}^{k}\sum_{j=i+1}^{k}E_{ij}\,\mathcal{P}'(x_i,x_j),\qquad \mathcal{P}'(x,y)=\mathcal{P}(x+y)-\mathcal{P}(x)-\mathcal{P}(y),\]
où les matrices publiques $E_{ij}\in\mathbb{F}_q^{m\times m}$ servent à multiplier par des puissances de $X$, modulo un polynôme irréductible de degré~$m$. Une signature est $S=(s_1,\dots,s_k)$ ; elle est valide si $\mathcal{P}^\star(S)=H(M)$. Ici $q=16$. Les paramètres $m$, $k$, $o$, $n$ définissent quatre jeux (MAYO-1, 2, 3 et 5), selon le niveau de sécurité demandé par le NIST~\cite{nistadd}. Les inconnues sont maintenant $k\cdot o$, ce qui permet d'avoir assez d'inconnues pour les $m$ équations.

\subsubsection{Évaluation par produits de matrices}
Les $m$ polynômes de $\mathcal{P}$ s'écrivent $p_\ell(x)=x^{\top}P_\ell\,x$, où $P_\ell$ est une matrice triangulaire supérieure. On la découpe en blocs :
\[P_\ell=\begin{pmatrix}P^{(1)}_\ell&P^{(2)}_\ell\\0&P^{(3)}_\ell\end{pmatrix},\]
où $P^{(1)}_\ell$ est de taille $v\times v$ avec $v=n-o$ (notée $V$ dans le code et dans le tableau~\ref{tab:params}), $P^{(2)}_\ell$ de taille $v\times o$ et $P^{(3)}_\ell$ de taille $o\times o$. Les blocs $P^{(1)}$ et $P^{(2)}$ sont recalculés à partir d'une graine par une fonction à base d'AES ; seul $P^{(3)}$ est stocké dans la clé publique~\cite{nibbling}. Soit $S$ la matrice dont les colonnes sont $s_1,\dots,s_k$. Tout se déduit de $S^{\top}P_\ell S$ : $\mathcal{P}(s_i)_\ell$ est l'élément $(i,i)$ de cette matrice, et $\mathcal{P}'(s_i,s_j)_\ell$ est la somme des éléments $(i,j)$ et $(j,i)$. \emph{Ce sont donc les produits de matrices qui coûtent le plus cher}, d'où l'intérêt de les optimiser~\cite{nibbling}.
\[P_\ell=\begin{pmatrix}P^{(1)}_\ell&P^{(2)}_\ell\\0&P^{(3)}_\ell\end{pmatrix},\]
où $P^{(1)}_\ell$ est de taille $v\times v$ avec $v=n-o$ (noté $V$ dans le code et au tableau~\ref{tab:params}), $P^{(2)}_\ell$ de taille $v\times o$ et $P^{(3)}_\ell$ de taille $o\times o$. Les blocs $P^{(1)}$ et $P^{(2)}$ sont dérivés d'une graine par une fonction d'expansion à base d'AES ; seul $P^{(3)}$ est stocké dans la clé publique~\cite{nibbling}. Si $S$ est la matrice dont les colonnes sont $s_1,\dots,s_k$, tout se déduit de $S^{\top}P_\ell S$ : $\mathcal{P}(s_i)_\ell$ est son élément diagonal $(i,i)$ et $\mathcal{P}'(s_i,s_j)_\ell$ la somme des éléments $(i,j)$ et $(j,i)$. \emph{Les produits matrice $\times$ matrice dominent donc le coût}, ce qui motive leur optimisation~\cite{nibbling}.

\subsubsection{Algorithmes (version simplifiée)}
\begin{encadre}
\textbf{KeyGen.} Tirer \code{seed_sk} ; en dériver \code{seed_pk} et $O\in\mathbb{F}_{16}^{v\times o}$ ; développer $\{P^{(1)}_\ell,P^{(2)}_\ell\}$ depuis \code{seed_pk} ; calculer $P^{(3)}_\ell=\mathrm{Upper}\big(-O^{\top}(P^{(1)}_\ell O-P^{(2)}_\ell)\big)$. Clé publique : \code{seed_pk} et $\{P^{(3)}_\ell\}$ ; clé secrète : \code{seed_sk}.\\[2pt]
\textbf{Sign}$(M)$. (1)~Tirer un sel et calculer $t=H(M\,\|\,\text{sel})\in\mathbb{F}_{16}^m$. (2)~Tirer $V\in\mathbb{F}_{16}^{k\times v}$ (les « vinaigres »). (3)~Pour chaque $\ell<m$ : $L_\ell=(P^{(1)}_\ell+P^{(1)\top}_\ell)\,O+P^{(2)}_\ell$ ; $M_\ell=V\,L_\ell$ ; $Y_\ell=V\,P^{(1)}_\ell\,V^{\top}$. (4)~Assembler le système linéaire $Ax=y$, où $A$ provient des $M_\ell$ et $y=t+(\text{combinaison des }Y_\ell)$. (5)~Si le système n'a pas de solution, retourner à l'étape~2 (nouveaux vinaigres) ; sinon tirer une solution $x$ au hasard. (6)~Signature : $S=(V+(OX)^{\top},X)$, où $X$ est $x$ mis en matrice $k\times o$.\\[2pt]
\textbf{Verify.} Redévelopper $P^{(1)},P^{(2)}$ ; calculer $t=H(M\,\|\,\text{sel})$ ; calculer $Y_\ell=S\begin{pmatrix}P^{(1)}_\ell&P^{(2)}_\ell\\0&P^{(3)}_\ell\end{pmatrix}S^{\top}$, combiner en $\mathcal{P}^\star(S)$ ; accepter si le résultat égale $t$.
\end{encadre}
La signature fait donc intervenir trois choses : des produits de matrices (pour $M_\ell$ et $Y_\ell$, à partir de la clé publique développée $P^{(1)}$, $P^{(2)}$), la résolution d'un système linéaire par élimination de Gauss en temps constant, et des reprises quand le système n'a pas de solution. D'après~\cite{nibbling}, ce sont le calcul des évaluations partielles de $\mathcal{P}$ et l'élimination de Gauss qui coûtent le plus ; l'assemblage de $A$ et de $y$ ne coûte presque rien.

Un exemple complet (génération de clé, signature, vérification) sur $\mathbb{F}_{16}$ avec de très petits paramètres est donné à l'annexe~\ref{app:jouet}.

\subsection{La signature dans le code de référence}\label{sec:code}
\subsubsection{Le déroulement d'une tentative}
Dans le code de référence (annexe~\ref{app:inventaire}), la fonction \code{mayo_sign_signature} calcule le sel et $t$, puis fait une boucle d'au plus 256 tentatives (un octet sert de compteur). Chaque tentative enchaîne les étapes suivantes (figure~\ref{fig:pipeline}) :
\begin{enumerate}[itemsep=1pt]
\item tirer et décoder les vinaigres $V$ (avec SHAKE256) ;
\item \code{m_calculate_M_and_VPV} : calculer les matrices $M_\ell$ et tous les produits $v_i P^{(1)} v_j$ (c'est ici que se trouvent les noyaux étudiés) ;
\item \code{compute_rhs} (second membre $y$) puis \code{compute_A} (matrice $A$) ;
\item \code{sample_solution} : résoudre $Ax=y$ ; en cas d'échec, nouvelle tentative ;
\item \code{m_compute_s} : assembler la signature $s_i=[\,v_i+Ox_i\mid x_i\,]$.
\end{enumerate}

\begin{figure}[ht]\centering
\begin{tikzpicture}[node distance=5mm,font=\small,box/.style={draw,rounded corners=2pt,fill=ice,minimum width=3.4cm,minimum height=0.7cm,align=center},key/.style={draw=alert,thick,rounded corners=2pt,fill=alert!12,minimum width=3.1cm,minimum height=0.7cm,align=center},>=Stealth]
 \node[box] (a) {sel, $t=H(M\|\text{sel})$};
 \node[box,right=8mm of a] (b) {tirer $V$ (SHAKE256)};
 \node[key,right=8mm of b] (c) {\code{m_calculate_M_and_VPV}\\[-1pt]{\scriptsize noyaux étudiés}};
 \node[box,below=14mm of c] (d) {\code{compute_rhs}, \code{compute_A}};
 \node[box,left=8mm of d] (e) {\code{sample_solution}\\[-1pt]{\scriptsize (Gauss, temps constant)}};
 \node[box,left=8mm of e] (f) {\code{m_compute_s}\\[-1pt]{\scriptsize signature}};
 \draw[->] (a)--(b); \draw[->] (b)--(c); \draw[->] (c)--(d); \draw[->] (d)--(e); \draw[->] (e)--(f);
 \draw[->,dashed] ($(e.north)+(0.9,0)$) -- node[right,font=\scriptsize,align=left]{échec :\\nouvelle\\tentative} ($(b.south)+(0.9,0)$);
\end{tikzpicture}
\caption{Une tentative de signature dans le code de référence. Les noyaux optimisés se trouvent dans l'étape en rouge, répétée à chaque tentative.}\label{fig:pipeline}
\end{figure}

\subsubsection{Où sont les noyaux étudiés}
Dans \code{m_calculate_M_and_VPV} (chemin à $L$ implicite, celui où les deux noyaux sont utilisés), les opérations se font dans cet ordre : construction des tables de multiplication de $V$ (\code{mayo_V_multabs}), puis :
\begin{itemize}[itemsep=1pt]
\item \code{P1_times_Vt} : accumule $P^{(1)}\cdot V^{\top}$ dans un tampon ; \code{Vt_times_Pv} en déduit $V P^{(1)}V^{\top}$, c'est-à-dire tous les $v_iP^{(1)}v_j$ ;
\item \code{P1t_times_Vt} : accumule la contribution transposée $P^{(1)\top}V^{\top}$ sur le même tampon, qui devient $(P^{(1)}+P^{(1)\top})V^{\top}$ (la diagonale s'annule en caractéristique~2) ; le calcul de $M_\ell$ réutilise ce produit, puis \code{Ut_times_O} et \code{Vt_times_L} complètent $M_\ell$.
\end{itemize}
Les noyaux \code{P1_times_Vt} et \code{P1t_times_Vt} calculent donc des produits de la \emph{partie vinaigre de la clé publique} $P^{(1)}$ par les multiples de $V$. Ce sont les seuls que ce mémoire optimise (variante A2), pour MAYO-1 et MAYO-5.

\subsubsection{Rôle de $P_1$ et paramètres}
$P_1$ est une matrice triangulaire supérieure dont chaque case est un $m$-vecteur de $\mathbb{F}_{16}^m$. Elle est rangée à plat, une entrée après l'autre : l'entrée numéro $e$ commence à la position $e\cdot b$. Chaque noyau calcule, pour chaque entrée de $P_1$ et chaque vecteur de $V$, le produit d'un nombre de $\mathbb{F}_{16}$ par le $m$-vecteur de l'entrée, et l'ajoute à un accumulateur de $L$ « limbes ». Les bornes des boucles ($V$, $K$) et les adresses de lecture ($e\cdot b$) ne dépendent que de paramètres publics. Le tableau~\ref{tab:params} donne les valeurs utiles.

\begin{table}[ht]\centering\small
\caption{Paramètres des noyaux étudiés. $m$ : longueur du $m$-vecteur ; $b=m/2$ : octets par $m$-vecteur ; $L$ : limbes de 64 bits d'un accumulateur ; $K$ : nombre de vecteurs de $V$ ; $V$ : dimension de la partie vinaigre ; $N_{\mathrm{ent}}=V(V+1)/2$ : nombre d'entrées de $P_1$.}\label{tab:params}
\begin{tabular}{@{}lrrrrrrr@{}}\toprule
Paramètre & $m$ & $b$ & $L$ & $K$ & $V$ & $N_{\mathrm{ent}}$ & $m \bmod 64$\\\midrule
MAYO-1 & 80 & 40 & 5 & 10 & 80 & 3\,240 & 16\\
MAYO-2 & 64 & 32 & 4 & -- & 73 & -- & 0\\
MAYO-3 & 108 & 54 & -- & -- & -- & -- & 44\\
MAYO-5 & 142 & 71 & 9 & 12 & 142 & 10\,153 & 14\\\bottomrule
\end{tabular}
\end{table}

\subsection{Représentation des données}\label{sec:repr}
\subsubsection{Le corps $\mathbb{F}_{16}$ et les $m$-vecteurs}
$\mathbb{F}_{16}$ est construit comme $\mathbb{F}_2[x]/(x^4+x+1)$ (commentaire du code, \code{simple_arithmetic.h}). Un élément tient dans un \emph{quartet} (4 bits) et l'addition est le OU exclusif. L'unité de calcul de MAYO est le \emph{$m$-vecteur} : la liste des éléments qui occupent la même case $(i,j)$ dans les $m$ matrices $P_\ell$. On calcule ainsi sur les $m$ composantes en parallèle, ce qui convient bien à la vectorisation~\cite{nibbling}.

\subsubsection{Empaquetage en quartets}
Un $m$-vecteur occupe $b=m/2$ octets : deux éléments par octet, l'élément d'indice pair dans les 4 bits de poids faible. Les accumulateurs sont découpés en \emph{limbes} de 64 bits, soit 16 éléments par limbe. Il faut donc $L=\lceil m/16\rceil$ limbes, avec éventuellement des quartets de bourrage (remplissage) dans le dernier limbe ($m=142$ : $L=9$, 144 positions, dont deux de bourrage). Cette représentation par quartets vient de~\cite{nibbling}, qui l'a préférée à la représentation « bitsliced » (par plans de bits) de la première version. Elle permet d'utiliser des instructions qui permutent les octets pour multiplier, au prix d'un changement du format de la clé.

\subsection{Arithmétique $\mathbb{F}_{16}$ par \vpshufb}\label{sec:arith}
\subsubsection{Rappel : le registre AVX2 et \vpshufb}
Un registre AVX2 \code{ymm} contient 256 bits, soit deux \emph{voies} de 128 bits (\code{xmm} désigne la moitié basse). La plupart des instructions, dont \vpshufb, travaillent séparément sur chaque voie. \vpshufb{} reçoit une table de 16 octets (une par voie) et un registre d'index. Pour chaque octet d'index, elle écrit l'élément de la table désigné par les 4 bits de poids faible de l'index, ou 0 si le bit de poids fort de l'index vaut 1. Un registre \code{ymm} fait donc 32 recherches avec une seule instruction.

\subsubsection{Principe}
Avec des éléments de 4 bits, multiplier un vecteur par un nombre $c$ revient à chercher chaque élément $x$ dans une table de 16 cases qui contient $c\cdot x$. L'instruction \vpshufb{} (\code{_mm256_shuffle_epi8}) fait cette recherche pour chaque octet d'un registre d'index. C'est la méthode des « quatre Russes »~\cite{russes} : on remplace des multiplications par des lectures dans une table. Pour calculer $A\cdot B$, on découpe $A$ en bandes de $t$ colonnes et $B$ en bandes de $t$ lignes. Pour chaque bande de $B$, on prépare toutes les combinaisons de ses lignes ($16^t$ entrées), puis chaque ligne de $A$ sert d'indice dans cette table. Ici la table tient dans un registre, donc $t=1$ (voir aussi~\cite{albrecht2011} pour $\mathbb{F}_2$). Les tables dépendent de $V$ (qui est secret), mais elles ne servent que de données pour \vpshufb{} : elles ne servent jamais d'adresse mémoire.

\subsubsection{Deux produits par recherche}
Pour une paire de nombres $(c_{2p},c_{2p+1})$, la table $T_p[x]=(c_{2p}x)\oplus\big((c_{2p+1}x)\ll4\big)$ donne \emph{les deux} produits dans un seul octet. Le code range $T_p$ dans un registre de 256 bits dont les deux voies sont identiques (lemme~\ref{lem:voies}). Un registre de données de 32 octets contient 64 éléments. Il faut séparer les quartets impairs et pairs (ET avec \code{0x0f}, puis décalage de 4 bits), donc deux \vpshufb{} par registre et par paire. Les accumulateurs « impairs » et « pairs » restent séparés jusqu'à la fin de la ligne, où on les remet ensemble. Cette technique est décrite dans l'article \emph{Nibbling MAYO}~\cite{nibbling}.

\subsubsection{Recouvrement en fin de vecteur}
Pour $m\le64$, $64<m\le128$ et $128<m\le192$, le code officiel utilise respectivement 1, 2 et 3 chargements de 32 octets. Le dernier est \emph{reculé} pour se terminer à l'octet $b$ (adresse \code{vec + M_VEC_LIMBS_MAX - 4} limbes). Les octets communs sont donc traités deux fois, et les écritures finales se recouvrent avec les mêmes valeurs. Pour MAYO-5, on a les chargements $[0,32)$, $[32,64)$ et $[40,72)$ : 192 éléments sont calculés pour 142 utiles, et le dernier chargement lit un octet de trop après l'entrée de $P_1$ (figure~\ref{fig:charge}).

\begin{figure}[ht]\centering
\begin{tikzpicture}[x=0.105cm,y=1cm,>=Stealth,font=\small]
 \def\h{0.45}
 % echelle : 1 unite = 1 octet, entree de 71 octets : 0..71
 \node[anchor=east] at (-2,4.2) {Entrée de $P_1$ (71 octets)};
 \fill[ice] (0,4) rectangle (71,4.5); \draw (0,4) rectangle (71,4.5);
 \foreach \x in {0,32,64,71}{\draw[muted] (\x,3.9)--(\x,4.6); \node[font=\scriptsize,muted] at (\x,4.85) {\x};}
 % officiel
 \node[anchor=east] at (-2,3) {Officiel : 3 ymm};
 \fill[steel!35] (0,2.8) rectangle (32,3.25); \draw (0,2.8) rectangle (32,3.25);
 \fill[steel!55] (32,2.8) rectangle (64,3.25); \draw (32,2.8) rectangle (64,3.25);
 \fill[alert!40] (40,2.25) rectangle (72,2.7); \draw (40,2.25) rectangle (72,2.7);
 \node[font=\scriptsize] at (16,3.02) {$[0,32)$}; \node[font=\scriptsize] at (48,3.02) {$[32,64)$}; \node[font=\scriptsize] at (56,2.47) {$[40,72)$ : recouvre $[40,64)$, lit 1 octet de trop};
 % A2
 \node[anchor=east] at (-2,1.2) {A2 SAFE : 2 ymm + 1 xmm};
 \fill[steel!35] (0,1) rectangle (32,1.45); \draw (0,1) rectangle (32,1.45);
 \fill[steel!55] (32,1) rectangle (64,1.45); \draw (32,1) rectangle (64,1.45);
 \fill[ok!40] (63,0.45) rectangle (71,0.9); \draw (63,0.45) rectangle (71,0.9);
 \node[font=\scriptsize] at (16,1.22) {$[0,32)$}; \node[font=\scriptsize] at (48,1.22) {$[32,64)$};
 \node[font=\scriptsize,anchor=east] at (62,0.67) {queue $[63,71)$, puis décalage de 8 bits};
\end{tikzpicture}
\caption{Chargements pour une entrée de $P_1$ de MAYO-5 (octets). Le noyau officiel recouvre l'intervalle $[40,64)$ et dépasse l'entrée d'un octet ; A2 SAFE lit exactement $[0,71)$.}\label{fig:charge}
\end{figure}



\subsection{Panorama des analyses menées}\label{sec:panorama}
Pour s'y retrouver dans la suite, voici les analyses et ce qu'elles établissent (détails aux chapitres~\ref{ch:methodo} et~\ref{ch:resultats}) :
\begin{itemize}[itemsep=1pt]
\item \textbf{Comptage} (théorique, dans un modèle précis) : au moins $\lceil m/32\rceil$ \vpshufb{} par paire de nombres et par $m$-vecteur ; A2 atteint cette limite pour $m\in\{80,142\}$ ;
\item \textbf{Correction et sûreté} : tests de réponses connues (KAT), tests unitaires, limites de lecture et d'écriture, propriétés des tables, indépendance vis-à-vis du bourrage ;
\item \textbf{Mesures de temps} sur une machine physique (compteur TSC, unité statistique : le processus) : gains sur les noyaux seuls, puis sur \code{mayo_sign}, avec contrôle de l'effet d'ordre et des sondes de mesure ;
\item \textbf{Analyse du code produit sans l'exécuter} (E1, E1b) : on compte les instructions et les écritures dans les boucles produites par GCC~16.2.0, et on désactive des passes du compilateur une à une ;
\item \textbf{Chaque résultat porte une étiquette} : \F{} fait établi, \I{} déduction, \Hy{} hypothèse non testée.
\end{itemize}

\section{Revue de littérature}\label{sec:revue}
\subsection{Implémentations existantes}
\begin{description}[leftmargin=0.5cm]
\item[Référence C.] Arithmétique ordinaire, sans \vpshufb{}.
\item[AVX2 officiel (MAYO-C).] Fichier \code{src/AVX2/shuffle_arithmetic.h} : macros \code{MULTIPLY_ACCUMULATE_M_VEC} et \code{DISENTANGLE_AND_ACCUMULATE_M_VECS}, avec le recouvrement décrit plus haut. C'est la base de ce travail : version amont \code{74de637}~\cite{mayoc}.
\item[\emph{Nibbling MAYO}~\cite{nibbling}.] Implémentation AVX2 et Cortex-M4 optimisée, avec la représentation en quartets. Au paragraphe 5.1, un \vpshufb{} donne 64 multiplications, les accumulateurs impairs et pairs sont séparés et, pour $m=96$ (MAYO-3 de l'époque), un recouvrement « duplique 32 opérations ». Le débit est d'un \vpshufb{} par cycle sur Skylake et de deux sur Ice Lake. Nous avons lu l'article en entier. Il utilise les paramètres du tour~1 ($m=64$ pour MAYO-1 et MAYO-2, $96$ pour MAYO-3, $128$ pour MAYO-5) : MAYO-1 et MAYO-5 n'y ont donc \emph{pas} de fin de vecteur à traiter. Les mesures ont été faites sur Skylake (Xeon E3-1245~v5) et Ice Lake (Xeon Gold 6338), avec \code{clang-14}, Turbo Boost désactivé, médiane de 10\,000 répétitions. Les tables de multiplication sont construites par la méthode de~\cite{bch23}, appliquée deux fois puis entrelacée. L'approche par \vpshufb{} se distingue de~\cite{bch23} en cherchant \emph{deux} produits par recherche. L'article évoque le recouvrement pour $m$ non multiple de 64 (cas MAYO-3), mais sans traitement spécial de la fin du vecteur. Le code publié avec l'article (\code{shuffle_arithmetic_\{64,96,128\}.h}) charge aux positions de limbes 0 et 2 (octets $[0,32)$ et $[16,48)$) et utilise quatre \vpshufb{} par paire pour $m=96$.
\item[Version GFNI.] Elle utilise l'instruction \code{vgf2p8affineqb} et lit sur toute la largeur dans une marge de sécurité (\code{MAYO_MVEC_SLACK}) ; elle est hors du cadre de ce mémoire.
\end{description}


\subsection{Fuites par canaux auxiliaires et temps constant}
Les attaques par mesure du temps~\cite{kocher} et par le cache~\cite{bernstein} exigent que les adresses mémoire et les branches du programme ne dépendent pas de données secrètes. Dans les noyaux, les tables de multiplication calculées à partir du vinaigre ne sont que des données de \vpshufb{} : elles ne servent ni d'adresses ni de conditions. A2 doit garder cette propriété ; la section~\ref{sec:surete} dit où en est la preuve.

\section{Évaluation et positionnement}
Le tableau~\ref{tab:etat} sépare ce qui était connu, ce qui est venu après l'article, et ce qu'A2 fait de différent. Il se fonde sur les sources consultées : le paragraphe 5.1 de l'article (lu en entier), les parties AVX2 du code, et tout l'historique du fichier \code{shuffle_arithmetic.h} dans le dépôt officiel (versions \code{3c9be9e}, \code{d81b689}, \code{11ff0f0}, \code{4b7cd94}, \code{fef05b5} et la dernière version de \code{main}, \code{74de637}).

\begin{table}[ht]\centering\small
\caption{Situation d'A2 par rapport aux sources consultées.}\label{tab:etat}
\begin{tabularx}{\linewidth}{@{}lXX@{}}\toprule
Catégorie & Élément & Source\\\midrule
Connu avant A2 & \vpshufb{} avec $t=1$, deux produits par octet ; séparation impair/pair ; remise en ordre à la fin & \cite{nibbling}\\
Connu avant A2 & Recouvrement de registres quand $m$ n'est pas multiple de 64 (MAYO-3, $m=96$) & \cite{nibbling}\\
Venu après l'article & Généralisation par chargements reculés pour tout $m\le192$ & code officiel (tour~2)\\
Différent dans A2 & Fin de vecteur $\le16$ éléments regroupée dans \emph{un} registre xmm (quartets [bas$\mid$haut]), avec \emph{un} \vpshufb{} xmm sur la voie basse de la table & ce travail\\
Différent dans A2 & Variante SAFE pour $m=142$ : on charge $[63,71)$ puis on décale, sans aucune lecture en dehors de l'entrée & ce travail\\
Différent dans A2 & Propriétés vérifiées pour la fin du vecteur (section~\ref{sec:surete}) & ce travail\\
Non revendiqué & Recouvrement, séparation impair/pair, remise en ordre tardive, schéma $t=1$, nouveauté absolue & --\\\bottomrule
\end{tabularx}
\end{table}

Nous n'avons trouvé aucun \code{vpshufb} xmm sur une fin de vecteur, ni dans le code publié avec l'article, ni dans l'historique officiel jusqu'à la dernière version de \code{main}. Les \code{__m128i} qui apparaissent à partir de \code{fef05b5} (lignes 548--606) servent à autre chose : transposer des $m$-vecteurs par décalage de bits, et charger une table de 16 octets répétée sur 256 bits. \textbf{Limite :} nous n'avons lu l'article en entier que pour sa section 5.1 et ses parties AVX2 utiles. Ne pas avoir trouvé la technique dans ces sources ne prouve pas qu'elle n'existe pas ailleurs.

